\TieudeT{PHONG TỎA VẬT LÍ}
\subsubsection*{PHẦN I. Thí sinh trả lời câu hỏi từ câu 1 đến câu 18. Mỗi câu hỏi thí sinh chỉ chọn một phương án}
\setcounter{ex}{0}
\begin{ex}%book
Công thức tổng quát để xác định vận tốc tổng hợp là
\choice
{$v_{13}=v_{12}+v_{23}$}
{\True $\overrightarrow{v}_{13}=\overrightarrow{v}_{12}+\overrightarrow{v}_{23}$}
{$\overrightarrow{v}_{13}=\overrightarrow{v}_{12}-\overrightarrow{v}_{23}$}
{$v_{13}=v_{12}-v_{23}$}
\loigiai{}
\haidong{7} 
\end{ex}
\begin{ex}%book
Vận tốc của vật so với hệ quy chiếu chuyển động được gọi là
\choice
{vận tốc kéo theo}
{\True vận tốc tương đối}
{vận tốc tuyệt đối}
{vận tốc tức thời}	
\loigiai{}
\haidong{7} \end{ex}

\begin{ex}%book
Vận tốc tổng hợp có giá trị bằng
\choice
{\True tổng các vận tốc thành phần}
{độ dịch chuyển thành phần}
{tổng các quãng đường đi}
{tổng các độ dịch chuyển}
\loigiai{
}
\haidong{7} \end{ex}	
\begin{ex}%book
Một hành khách ngồi trong xe A, nhìn qua cửa sổ thấy xe B bên cạnh và sân ga đều chuyển động như nhau. Phát biểu nào sau đây đúng?
\choice
{\True Xe A chuyển động so với xe B, xe B đứng yên so với sân ga}
{Xe A đứng yên so với xe B, xe B chuyển động so với sân ga}
{Xe A và xe B đều đứng yên so với sân ga}
{Xe A và xe B đều chuyển động so với sân ga}

%\loigiai{
%}
\haidong{7} \end{ex}	
\begin{ex}%book %[1P1K9-1][Loai1] 
Một chiếc máy bay đang bay từ Thành phố Hồ Chí Minh đến Thủ đô Hà Nội với tốc độ $525\ km/h$. Trong ngày hôm đó, gió thổi về hướng Nam với tốc độ $36\ km/h$. Xem như máy bay chuyển động thẳng đều theo hướng Bắc và quãng đường bay từ Thành phố Hồ Chí Minh đến Thủ đô Hà Nội là $1160\ km$. Thời gian bay của máy bay trên quãng đường đó là
\choice
{\True $2,37h$}
{$2,07h$}
{$32,2h$}
{$2,2h$}
\haidong{13}
\loigiai{		
\begin{itemize}
\item Gọi vận tốc của máy bay và vận tốc của gió đối với mặt đất lần lượt là $\vec v_{12}$ và $\vec v_{23}$, vận tốc tương đối của máy bay đối với gió là $\vec v_{13}$.
\item Công thức cộng vận tốc: $\vec v_{13}=\vec v_{12}+\vec v_{23}$
\item Theo bài ra hướng gió và hướng chuyển động của máy bay ngược chiều nhau nên: $\vec v_{12}\uparrow \downarrow \vec v_{23}$.			
\item Chọn chiều dương là chiều chuyển động của máy bay: $v_{13}=525-36=489\ km/h$
\item Thời gian bay của máy bay: $t=\dfrac{d}{v_{13}}=\dfrac{1160}{489}=2,37\ h$.
\end{itemize}	
}
\end{ex}	
\begin{ex}%book %[1P1G9-1][Loai1]
Một xe đạp chuyển động thẳng đều với vận tốc lúc không có gió là $15\ km/h$. Người này đi từ $A$ tới $B$ xuôi gió và đi từ $B$ trở lại $A$ ngược gió. Vận tốc gió là $1\ km/h$. Khoảng cách $AB = 28\ km$. Thời gian tổng cộng cả đi và về là
\choice
{$1,25\ h$}
{$2,5\ h$}
{$1,75\ h$}
{\True $3,75\ h$}	
\haidong{13}
\loigiai{
\begin{itemize}
\item Vận tốc xe khi xuôi gió: $v=15+1=16\ km/h$.
\item Thời gian đi xuôi gió: $t_1=\dfrac{28}{16}=1,75\ h$.
\item Vận tốc xe khi ngược gió: $v=15-1=14\ km/h$.
\item Thời gian xe đi ngược gió: $t_2=\dfrac{28}{14}=2\ h$.
\item Thời gian tổng cộng đi và về là: $t=t_1+t_2=3,75\ h$.
\end{itemize}	
}
\end{ex}	
\begin{ex}%book %[1P1-9.1K][Loai1] 
Hai đoàn tàu $A$ và $B$ chạy song song ngược chiều nhau. Đoàn tàu $A$ dài $150\ m$ chạy đều với vận tốc $15\ m/s$. Đoàn tàu $B$ chạy đều với vận tốc $10\ m/s$. Một hành khách đứng bên cửa sổ của tàu $B$ sẽ nhìn thấy đoàn tàu $A$ qua trước mặt mình trong
\choice
{$10\ s$}
{$15\ s$}
{$30\ s$}
{\True $6\ s$}
\haidong{13}
\loigiai{
\begin{itemize}
\item Độ lớn vận tốc của tàu B so với tàu A là: $v_{B/A} = v_B + v_A = 10 + 15 = 25$ m/s.
\item Thời gian mà tàu A qua trước mặt một hành khác trên tàu B là: $t = \dfrac{d}{v_{B/A}} = \dfrac{150}{25} = 6\ s$.
\end{itemize}	}
\end{ex}	
\begin{ex}%book %[1P1K9-2][Loai1]
Ô tô $A$ chạy thẳng đều theo hướng Tây với vận tốc $40\ km/h$. Ô tô $B$ chạy thẳng đều về hướng Bắc với vận tốc $60\ km/h$. Vận tốc của ô tô $B$ đối với người ngồi trên ô tô $A$ là
\choice
{\True $72,11\ km/h$}
{$56,23\ km/h$}
{$65,56\ km/h$}
{$78,21\ km/h$}	
\haidong{13}	
\loigiai{
\begin{itemize}
\item Áp dụng công thức cộng vận tốc ta có: $\overrightarrow{v}_{AB}=\overrightarrow{v}_{AD}+\overrightarrow{v}_{DB}$  
\item Ta có hai xe đi theo hai hướng vuông góc nhau $\Rightarrow v_{AB}=\sqrt{v_{AD}^2+v_{BD}^2}=\sqrt{40^2+60^2}=20\sqrt{13}\approx72,11\ km/h$.
\end{itemize}	
}
\end{ex}


\begin{ex}%book %[1P1-9.2G][Loai1]
Một người lái xuồng máy dự định mở máy cho xuồng chạy ngang sông (vuông góc với dòng chảy). Nhưng do nước chảy nên khi sang đến bờ bên kia, thuyền cách địa điểm của bến dự định là $180\ m$ về phía hạ lưu và mất $1$ phút. Biết chiều rộng của sông là $240\ m$. Vận tốc của xuồng so với bờ sông là 
\choice
{$6\ m/s$}
{$3\ m/s$}
{$4\ m/s$}
{\True $5\ m/s$}
\haidong{13}	
\loigiai{
\immini[thm]{\begin{itemize}
\item Gọi 1 là xuồng, 2 là nước, 3 là bờ, thì:
\begin{itemize}
\item Vận tốc của xuồng so với nước là $\vec{v}_{12}$
\item Vận tốc của nước so với bờ là $\vec{v}_{23}$
\item Vận tốc của xuồng so với bờ là $\vec{v}_{13}$
\end{itemize}
\item Do mũi xuồng vuông góc với dòng nước và xuồng trôi đến C nên các vector vận tốc $\vec{v}_{12}$ và $\vec{v}_{13}$ được biểu diễn như hình vẽ.
\item Vận tốc của xuồng đối với nước sông: $v_{12}=\dfrac{AB}{t}=\dfrac{240}{60}=4\ m/s$
\item Vận tốc của nước so với bờ là: $v_{23}=\dfrac{BC}{t}=\dfrac{180}{60}=3\ m/s$
\item Áp dụng công thức cộng vận tốc: $\vec{v}_{13}=\vec{v}_{12}+\vec{v}_{23}$	(*)
\item Bình phương hai vế (*) ta có: $v_{13}=\sqrt{v_{12}^2+v_{23}^2}=5\ m/s$
\end{itemize}}{\begin{tikzpicture}[scale=1,thick,>=stealth']
\tkzDefPoints{0/0/m, 4/0/n, 0/3/p, 4/3/q, 1.5/0/a, 1.5/3/b, 3.5/3/c, 1.5/3.3/b1, 3.5/3.3/c1}
\path (a) -- ($(a)!2cm!0:(c)$)coordinate(d);
\tkzDefPointBy[projection=onto a--b](d) \tkzGetPoint{a1}
\tkzDefPointBy[projection=onto a--n](d) \tkzGetPoint{a2}
\tkzDrawSegments(m,n p,q)
\tkzDrawSegments[dashed](a1,d a2,d a,b a,c)
\draw[very thick,red,->](a)--(d)node[below right,black]{$\vec{v}_{13}$};
\draw[very thick,blue,->](a)--(a1)node[left,black]{$\vec{v}_{12}$};
\draw[very thick,blue,->](a)--(a2)node[below,black]{$\vec{v}_{23}$};
\tkzDrawPoints[size=3](a,b,c)
\tkzMarkRightAngles(a,b,c)
\draw[dashed,<->] (b1)--node[midway,above]{$180\ m$}(c1);
\draw[dashed,<->] (m)--node[midway,above,rotate=90]{$240\ m$}(p);
\node at (b)[above left]{$B$};
\node at (a)[below]{$A$};
\node at (c)[above right]{$C$};
\end{tikzpicture}}
}
\end{ex}	
\begin{ex}%Tailieuchuan
Một chiếc thuyền chuyển động ngược dòng với tốc độ $14,4 \ km/h$ so với mặt nước, nước chảy với tốc độ $2 \ m/s$ so với bờ. Tốc độ của thuyền so với bờ là
\choice
{\True $2 \ m/s$}
{$6 \ m/s$}
{$12,4 \ km/h$}
{$16,4 \ km/h$}
\loigiai{

}
\end{ex}
\begin{ex}%Tailieuchuan
Một ô tô chạy theo hướng Đông với vận tốc $v_1$, một xe máy chạy thẳng về hướng Bắc với vận tốc $v_2$. Vận tốc của người chạy xe máy đối với người ngồi trên ô tô được xác định theo biểu thức nào sau đây?
\choice
{$v_2-v_1$}
{$v_2+v_1$}
{\True $\sqrt{v_1^2+v_2^2}$}
{$0$}
\loigiai{

}
\end{ex}
\begin{ex}%Tailieuchuan
Một đoàn tàu đi theo hướng Bắc - Nam với vận tốc $v_1$ so với mặt đường và trên tàu có một hành khách đi về phía đầu tàu với vận tốc $v_2$ so với mặt sàn của tàu. Vận tốc của người đó so với mặt đường được xác định thông qua biểu thức nào sau đây?
\choice
{\True $v_1+v_2$}
{$v_1-v_2$}
{$\sqrt{v_1^2-v_2^2}$}
{$\sqrt{v_1^2+v_2^2}$}
\loigiai{

}
\end{ex}
\begin{ex}%Tailieuchuan
Hai xe ô tô chạy ngược chiều trên một đoạn đường thẳng với vận tốc của xe thứ nhất và xe thứ hai có độ lớn lần lượt là $120\ km/h$ và $80\ km/h$. Chọn chiều dương ngược chiều chuyển động của xe thứ nhất. Vận tốc của xe thứ nhất so với xe thứ hai có giá trị đại số là
\choice
{$40\ km/h$}
{$200\ km/h$}
{$-40\ km/h$}
{\True $-200\ km/h$}
\loigiai{

}
\end{ex}
\begin{ex}%Tailieuchuan
Một chiếc ghe đi xuôi dòng từ bến đến chợ rồi ngược dòng quay về. Cho biết vận tốc của ghe so với nước là $15 \ km/h$, vận tốc của nước so với bờ là $3 \ km/h$. Biết quãng đường từ bến đến chợ dài $18\ km$. Thời gian chuyển động của ghe là
\choice
{$2$ giờ}
{\True $2,5$ giờ}
{$3$ giờ}
{$3,5$ giờ}
\loigiai{

}
\end{ex}
\begin{ex}%Tailieuchuan
Hai bên sông $A B$ cách nhau $70\ km$, thời gian một ca nô khi xuôi dòng $A$ đến $B$ ít hơn $48$ phút so với thời gian ca nô khi ngược dòng từ $B$ đến $A$. Tốc độ của ca nô trong nước yên lặng là $30 \ km/h$. Vận tốc của dòng nước là
\choice
{$5 \ m/s$}
{\True $5 \ km/h$}
{$28 \ m/s$}
{$28 \ km/h$}
\loigiai{}
\end{ex}
\begin{ex}%Tailieuchuan
Một hành khách ngồi trên toa tàu lửa số $1$ đang chuyển động với vận tốc $15\ m/s$ quan sát qua khe cửa thấy một đoàn tàu số $2$ chạy cùng chiều trên đường sắt bên cạnh (biết tàu số $1$ chạy nhanh hơn tàu số $2$). Từ lúc nhìn thấy điểm cuối đến lúc nhìn thấy điểm đầu của đoàn tàu số $2$ mất hết $8\ s$, biết tàu gồm $20$ toa, mỗi toa dài $4\ m$. Tốc độ của đoàn tàu số $2$ là
\choice
{\True $5 \ m/s$}
{$7 \ m/s$}
{$10 \ m/s$}
{$12 \ m/s$}
\loigiai{}
\end{ex}
\textit{Sử dụng thông tin sau cho Câu 17 và Câu 18:} Một người điều khiển thiết bị bay cá nhân bay theo hướng từ $A$ đến $B$. Gió thổi với vận tốc không đổi $27\ km/h$ theo hướng Bắc. Hướng AB lệch với hướng Bắc $60^\circ$ về phía Đông. 
\begin{ex}%book %[1P1K9-1][Loai1] 
Để bay theo đúng hướng từ $A$ đến $B$, với vận tốc tổng hợp là $54\  km/h$, người lái phải hướng thiết bị theo hướng
\choice
{\True Đông}
{Tây}
{Nam}
{Bắc}
\loigiai{}
\end{ex}

\begin{ex}
Bay được $6\ km$, thiết bị quay đầu bay về $A$ với vận tốc tổng hợp có độ lớn là $45\ km/h$ đúng hướng $B$ đến $A$. Tốc độ trung bình của thiết bị trên cả quãng đường bay là
\choice
{$46,8\ km/h$}
{\True $40,7\ km/h$}
{$54\ km/h$}
{$45\ km/h$}
\loigiai{
\begin{itemize}
\item  Gọi $\vec v_1$ là vận tốc của gió, có phương thẳng đứng, hướng về phía bắc và có độ lớn bằng 27 km/h.
\item Gọi $\vec v$ là vận tốc tổng hợp của thiết bị theo hướng A đến B và có độ lớn bằng 54 km/h.
\item Gọi $\vec v_2$ là vận tốc cần thiết của thiết bị để khi người điều khiển thiết bị đó có thể bay đúng hướng A đến B.
\item Quãng đường thiết bị bay từ A là 6 km, tốc độ bay tổng hợp là 54 km/h nên tốc độ bay của thiết bị là:\\
$v_2=\sqrt{v^2-v_1^2}=\sqrt{54^2-27^2}=27\sqrt{3}\ km/h$.
\item Vậy thời gian bay từ A là: $t_1=\dfrac{s_1}{v_2}=\dfrac{6}{27\sqrt{3}}\approx 0,128\ h$.
\item Quãng đường thiết bị bay về A là 6 km, tốc độ bay là 45 km/h nên tốc độ bay của thiết bị là:\\ $v'_2=\sqrt{v'^2-v_1^2}=\sqrt{45^2-27^2}=36\ km/h$.
\item Thời gian bay về A là: $t_2=\dfrac{s_2}{v'_2}=\dfrac{6}{36}\approx 0,167\ h$
\item Tốc độ trung bình của thiết bị trên cả quãng đường bay:\\ $v_{tb}=\dfrac{s}{t}=\dfrac{s_1+s_2}{t_1+t_2}=\dfrac{6+6}{0,128+0,167}\approx 40,68\ km/h$.
\end{itemize}
}
\end{ex}


































\subsubsection*{PHẦN 2. Thí sinh trả lời từ câu 1 đến câu 4. Trong mỗi ý a), b), c), d) ở mỗi câu, thí sinh chọn đúng hoặc sai.}
\setcounter{ex}{0}
\begin{ex}
\immini[thm]{Đứt gãy đá là hiện tượng đá bị tách ra và các mặt đối diện trượt đối với nhau phân chia thành lớp đá $A$ và lớp đá $B$. Trước khi lớp đá ở trên mặt trước trượt xuống sang phải, hai điểm $AB$ trùng nhau. Độ dịch chuyển toàn phần $\vec{AB}$ nằm dọc theo mặt phẳng đứt. Biết rằng $AD=17,0 \ m$; $AC=22,0 \ m$. 
}{\includegraphics[width=4cm]{images/2025_06_04_af03f1d88c157fe9bf98g-029}}
\choiceTF[t]
{\True Độ cao lớp đá $B$ có độ cao thấp hơn so với lớp đá $A$}
{Độ dịch chuyển toàn phần $\vec{AB}$ có độ lớn là $39,0 \ m$}
{Độ dịch chuyển toàn phần $\vec{AB}$ lệch so với mặt trên của tấm đá $A$ là $73,5^\circ$}
{Nếu mặt phẳng đứt nghiêng $52^\circ$ so với mặt phẳng ngang thì điểm $B$ dịch chuyển xuống theo phương thẳng đứng một đoạn $10,5 \ m$}
\loigiai{
\begin{itemchoice}
\itemch Vì lớp đá B trượt xuống so với lớp đá A nên lớp đá B thấp hơn lớp đá A.
$\Rightarrow$ Đúng
\itemch Độ dịch chuyển toàn phần của AB là: $AB=\sqrt{AD^2+AC^2}=\sqrt{17,0^2+22,0^2}=77,3 \ m$.
$\Rightarrow$ Sai
\itemch $\cos \widehat{BCA}=\dfrac{AC}{AB}=\dfrac{22,0}{77,3} \rightarrow \widehat{BCA}=73,5^{0}$
$\Rightarrow$ Đúng
\itemch Độ dịch chuyển theo phương thẳng đứng của lớp đá B là
$$
d_{dB}=AD \cdot \cos 52^\circ=17,0 \cdot \cos 52^\circ=10,5 \ m
$$

$\Rightarrow$ Đúng
\end{itemchoice}
}
\end{ex}
\begin{ex}
Một người bơi từ bờ này sang bờ kia của một con sông rộng $20\ m$ theo hướng vuông góc với bờ sông. Do nước sông chảy mạnh nên quãng đường người đó bơi gấp $2$ lần so với chiều rộng con sông.
\choiceTF[t]
{Người này sẽ bị nước đẩy dịch về hướng thượng nguồn}
{\True Độ dịch chuyển và quãng đường của người này có độ lớn bằng nhau và bằng $40\ m$ khi bơi sang bờ sông bên kia}
{Hướng dịch chuyển của người này làm với bờ sông một góc $60^\circ$}
{Vị trí điểm tới cách điểm đối diện với điểm khởi hành của người bơi là $60\ m$}
\loigiai{
\begin{itemchoice}
\itemch
\itemch
\itemch
\itemch	
\end{itemchoice}
}
\end{ex}
\begin{ex}
Một chiếc thuyền đi từ $A$ đến $B$ trên dòng sông cách nhau $30\ km$ rồi quay trở về $A$, nước chảy trên sông theo hướng từ $A$ đến $B$ với vận tốc $2\ km/h$. Biết thuyền đi với tốc độ $8\ km/h$ khi nuớc đứng yên. Chọn chiều dương từ $A$ đến $B$.
\choiceTF[t]
{\True Vận tốc của thuyền khi xuôi theo dòng nước có độ lớn là $10\ km/h$}
{\True Vận tốc của thuyền khi đi ngược dòng nước có độ lớn là $6 \ km / h$}
{Thời gian cả đi lẫn về của thuyền là $6$ giờ}
{Tốc độ trung bình của thuyền trong cả quá trình là $10 \ km / h$}
\loigiai{
\begin{itemchoice}
\itemch 
Gọi $\vec{v}_{12}$ là vận tốc của thuyền so với nước; $\vec{v}_{23}$ là vận tốc của nước so vói bờ; $\vec{v}_{13}$ là vận tốc của thuyền so với bờ. Chọn chiều dương từ A đến B. Áp dụng công thức cộng vận tốc: $\vec{v}_{13}=\vec{v}_{12}+\vec{v}_{23}$

Vận tốc của thuyền khi xuôi theo dòng nước: $v_{13}=v_{12}+v_{23}=8+2=10 \ km / h$.
$\Rightarrow$ Đúng
\itemch Vận tốc của thuyền khi đi ngược dòng nước: $v_{13}^{\prime}=v_{12}-v_{23}=8-2=6 \ km / h$.
$\Rightarrow$ Đúng
\itemch Thời gian thuyền xuôi dòng là $\dfrac{s}{v_{13}}=\dfrac{30}{10}=3 \ h$; Thời gian thuyền ngược dòng là $t_2=\dfrac{s}{v_{13}^{\prime}}=\dfrac{30}{6}=5 \ h$

Thời gian cả đi lẫn về của thuyền là: $t=t_1+t_2=8 \ h$.
$\Rightarrow$ Sai
\itemch Tốc độ trung bình của thuyền trong cả quá trình là $v=\dfrac{s_1+s_2}{t}=\dfrac{30+30}{8}=7,5 \ km / h$.
$\Rightarrow$ Sai
\end{itemchoice}
}
\end{ex}
\begin{ex}
Một đoàn tàu đang chuyển động thẳng đều với vận tốc $8\ m/s$, người soát vé đang ổn định hành khách trong toa tàu. Một học sinh đứng bên đường quan sát. Chọn chiều dương là chiều chuyển động của đoàn tàu.
\choiceTF[t]
{\True Vận tốc của tàu so với em học sinh có độ lớn là $8 \ m/s$}
{\True Khi người soát vé đứng yên trên tàu thì vận tốc của người soát vé đối với em học sinh trên có độ lớn là $8 \ m/s$}
{Khi người soát vé đi với vận tốc $1 \ m/s$ về phía đuôi tàu thì vận tốc của người soát vé đối với em học sinh trên có độ lớn là $9 \ m/s$}
{Khi người soát vé đi với tốc độ $1 \ m/s$ về phía đầu tàu thì vận tốc của người soát vé đối với em học sinh trên có độ lớn là $7 \ m/s$}
\loigiai{
\begin{itemchoice}
\itemch Do em học sinh đứng yên nên vận tốc của đoàn tàu so với mặt đường bằng vận tốc của đoàn tàu so với em học sinh và bằng $8 \ m/s$.
\itemch
Gọi $\vec{V_{12}}$ là vận tốc của người soát vé so với đoàn tàu.

$\vec{v_{23}}$ là vận tốc của đoàn tàu so với học $\sinh \Rightarrow v_{23}=8 \ m/s$.

$\vec{V_{13}}$ là vận tốc của người soát vé so với học sinh.

Áp dụng công thức cộng vận tốc: $\vec{v_{13}}=\vec{v_{12}}+\vec{v_{23}}$

Người soát vé đứng yên trên tàu nên $v_{12}=0$.

Vận tốc của người soát vé đối với học sinh là $v_{13}=v_{12}+v_{23}=8 \ m/s$
\itemch Khi người soát vé đi với tốc độ $1 \ m/s$ về phía đuôi tàu thì $v_{12}=-1 \ m/s$.\\
 Vận tốc của người soát vé đối với học sinh là 

$v_{13}=v_{12}+v_{23}=-1+8=7 \ m/s$.
\itemch Khi người soát vé đi với tốc độ $1 \ m/s$ về phía đầu tàu thì $v_{12}=1 \ m/s$.

Vận tốc của người soát vé đối với học $\sinh$ là $v_{13}=v_{12}+v_{23}=1+8=9 \ m/s$.
\end{itemchoice}
}
\end{ex}

\subsubsection*{PHẦN 3. Thí sinh trả lời từ câu 1 đến câu 6.}
\setcounter{ex}{0}
\textit{Sử dụng thông tin sau cho Câu 1 và Câu 2:} Một chiếc thuyền buồm chạy thẳng đều ngược dòng sông, sau $1$ giờ đi được $10\ km$. Một khúc gỗ trôi đều theo dòng sông, sau $1$ phút trôi được $\dfrac{100}{3}\ m$. 
\begin{ex}%book %[1P1K9-1][Loai1] 
Vận tốc của thuyền buồm so với bờ là bao nhiêu $km/h$?
\shortans[oly]{10}
\loigiai{
\begin{itemize}
\item Đổi $t_1 = 1h = 3600\ s_1 = 10\ km = 10000\ m,\ t_2 = 1\ p = 60\ s,\ s_2=\dfrac{100}{3}\ m.$
\item Gọi thuyền: 1; nước 2; bờ: 3
\item Vận tốc của thuyền so với bờ có độ lớn là:	$|v_{13}|=\dfrac{s_1}{t_1}=\dfrac{25}{9}\ m/s=10\ km/h$.	
\end{itemize}
}
\end{ex}	
\begin{ex}
Vận tốc của thuyền buồm so với nước là bao nhiêu $km/h$?
\shortans[oly]{12}
\loigiai{		
\begin{itemize}
\item Vận tốc của thuyền buồm so với bờ có độ lớn 10 km/h và chuyển động ngược chiều dòng nước.
\item[b.] Vận tốc của nước so với bờ có độ lớn là:
$|v_{23}|=\dfrac{s_2}{t_2}=\dfrac{5}{9}\ m/s$.
\item Áp dụng công thức cộng vận tốc: $\vec v_{12} = \vec v_{13} + \vec v_{32}\ \  hay\ \  \vec v_{12} = \vec v_{13} - \vec v_{23}$.
\item Chọn chiều dương là chiều chảy của dòng nước. 
\item Vì thuyền chảy ngược dòng nước nên $\vec v_{13}$ hướng ngược chiều dương, $\vec v_{23}$ hướng theo chiều dương.
\item Vận tốc của thuyền buồm so với nước có độ lớn 12 km/h và chuyển động ngược chiều dòng nước.
\end{itemize}}
\end{ex}
\begin{ex}%Tailieuchuan
Một chiếc thuyền chuyển động thẳng xuôi dòng nước từ bến $A$ tới bến $B$ cách nhau $6\ km$ dọc theo dòng sông rồi quay về $B$ mất $2$ giờ $30$ phút. Biết rằng tốc độ của thuyền trong nước im lặng là $5\ km/h$. Thuyền đi xuôi dòng mất thời gian bao nhiêu giờ?
\shortans[oly]{1}
\loigiai{
}
\end{ex}
\begin{ex}%Tailieuchuan
Một thủy thủ lái phà vận chuyển xe ô tô xuôi dòng với tốc độ $40\ km/h$ so với bờ. Sau khi hoàn thành công việc, thủy thủ lái phà quay lại theo lộ trình cũ với tốc độ $30\ km/h$ so với bờ. Biết rằng chiều và tốc độ của dòng nước không đổi trong suốt quá trình phà di chuyển và tốc độ của phà so với nước là không đổi. Tốc độ của dòng nước so với bờ là bao nhiêu $km/h$ (làm tròn kết quả đến chữ số hàng đơn vị)?
\shortans[oly]{5}
\loigiai{
}
\end{ex}
\begin{ex}%Tailieuchuan
Một người bơi trên một dòng sông yên lặng với tốc độ $2,5\ m/s$. Tốc độ tổng hợp của người đó khi bơi theo hướng Đông sang Tây và nước chảy tốc độ $1\ m/s$ theo hướng Bắc - Nam là bao nhiêu $m/s$ (làm tròn kết quả đến chữ số hàng phần trăm)?
\shortans[oly]{2,69}
\loigiai{
}
\end{ex}
\begin{ex}%Tailieuchuan
Một người đang ở phía Tây của một cái hồ và muốn bơi ngang qua đến vị trí ở phía Đông, đối diện với vị trí xuất phát của mình. Người này có thể bơi với vận tốc $1,9\ m/s$ khi nước yên lặng. Biết rằng lá cây trôi trên mặt nước hồ được $4,2\ m$ về hướng Nam trong $5\ s$. Nếu hồ rộng $4,8\ km$ thì người đó phải bơi trong bao nhiêu giây (làm tròn kết quả đến chữ số hàng đơn vị)?
\shortans[oly]{2823}
\loigiai{
}
\end{ex}









\newpage 
\TieudeT{PHONG TỎA VẬT LÍ}
\subsubsection*{PHẦN I. Thí sinh trả lời câu hỏi từ câu 1 đến câu 18. Mỗi câu hỏi thí sinh chỉ chọn một phương án}
\setcounter{ex}{0}
%mac
%\loai{Công thức}

\begin{ex}%book
Hai vận tốc cùng hướng thì tổng hợp hai vận tốc đó tính theo công thức nào sau đây?
\choice
{$v_{1,3}=v_{1,2}-v_{2,3}$}
{\True $v_{1,3}=v_{1,2}+v_{2,3}$}
{$v_{1,3}=v_{2,1}+v_{3,2}$}
{$v_{1,3}=v_{1,2}+v_{3,2}$}

%\loigiai{
%}
\haidong{7} \end{ex}
%\loai{Xác định loại vận tốc}


\begin{ex}%book
Vận tốc của vật đối với với hệ quy chiếu đứng yên được gọi là
\choice
{\True vận tốc tuyệt đối}
{vận tốc toàn phần}
{vận tốc kéo theo}
{vận tốc tương đối}

%\loigiai{
%}
\haidong{7} \end{ex}


%\loai{Xác định chuyển động, đứng yên}
\begin{ex}%book
Người lái đò đang ngồi yên trên chiếc thuyền thả trôi theo dòng nước thì
\choice
{người lái đò chuyển động so với dòng nước}
{\True người lái đò đứng yên so với dòng nước}
{người lái đò đứng yên so với bờ sông}
{người lái đò chuyển động so với chiếc thuyền}
%\loigiai{
%}
\haidong{7} \end{ex}

\begin{ex}%book
Một ô tô chở khách đang chạy trên đường. Đối với người nào dưới đây ô tô đang đứng yên?
\choice
{Người lái ô tô con đang vượt qua xe khách}
{Người đang đứng bên lề đường}
{Một người lái xe máy đang bị xe khách vượt qua}
{\True Một hành khách đang ngồi trong xe}
%\loigiai{
%}
\haidong{7} \end{ex}
%\loai{Nhiều chọn lựa}

%\loai{Đại lượng tương đối}
\begin{ex}%book
Phát biểu nào sau đây là \textbf{không} đúng? Đại lượng có tính tương đối là
\choice
{\True khối lượng}
{độ dịch chuyển}
{vận tốc}
{quỹ đạo}
%\loigiai{
%}
\haidong{7} \end{ex}

%mac
%\loai{Tính vận tốc kéo theo}
\begin{ex}%book %[1P1K9-1][Loai1] 
Một máy bay chuyển động thẳng đều từ $A$ đến điểm $B$ cách nhau $900\ km$ theo chiều gió mất $2$ giờ $30$ phút. Vận tốc của máy bay khi không có gió là $300\ km/h$. Vận tốc của gió là
\choice
{$360\ km/h$}
{\True $60\ km/h$}
{$420\ km/h$}
{$180\ km/h$}
\haidong{13}
\loigiai{		
\begin{itemize}
\item Đổi 2h30 phút $=$ 2,5 h.
\item Khi có gió thì vận tốc của máy bay so với gió là $v_{mb/g} = \dfrac{900}{2,5} = 360\ km/h$.
\item Theo công thức cộng vận tốc $v_{mb/g} = v_{mb/md} + v_{g/md} \Rightarrow v_{g/md} = v_{mb/g} - v_{mb/md} = 360 - 300 = 60$ km/h.
\end{itemize}	}
\end{ex}


\begin{ex}%book %[1P1K9-1][Loai1]
Một chiếc thuyền chuyển động thẳng đều với vận tốc $5\ km/h$ so với bờ. Một em bé đi đều từ đầu thuyền tới cuối thuyền với vận tốc $6\ km/h$ so với thuyền. Vận tốc của em bé so với bờ là
\choice
{\True $1\ km/h$}
{$2\ km/h$}
{$6\ km/h$}
{$5\ km/h$}
\haidong{13}
\loigiai{
\begin{itemize}
\item Ta có: $v_{TB}=5\ km/h;\ v_{ET}=1\ km/h$.
\item Áp dụng công thức cộng vận tốc ta có: $\overrightarrow{v}_{EB}=\overrightarrow{v}_{ET}+\overrightarrow{v}_{TB} \Rightarrow v_{EB}=v_{ET}-v_{TB}=6-5=1\ km/h$.
\end{itemize}	
}
\end{ex}
%\loai{Tự luận}

%\loai{Tính thời gian đi xuôi ngược}
\begin{ex}%book %[1P1-9.1G][Loai1]
Hai bến sông $A$ và $B$ cách nhau $18\ km$ theo đường thẳng. Biết rằng vận tốc của canô khi nước không chảy là $16,2\ km/h$ và vận tốc của dòng nước so với bờ là $1,5\ m/s$. Thời gian để canô đi từ $A$ đến $B$ rồi trở lại ngay từ $B$ tới $A$ là
\choice
{$1$ giờ $40$ phút}
{$30$ phút}
{$50$ phút}
{\True $2$ giờ $30$ phút}
\haidong{13}
\loigiai{
\begin{itemize}
\item Ta có: $v_{CN}=16,2\ km/h;\ v_{NB}=1,5\ m/s=5,4\ km/h$.
\item Áp dụng công thức cộng vận tốc ta có: $\overrightarrow{v}_{CB}=\overrightarrow{v}_{CN}+\overrightarrow{v}_{NB}$.
\item Khi xuôi dòng thì: $v_{CB}=v_{CN}+v_{NB}=16,2+5,4=21,6\ km/h$.
\item Thời gian canô đi từ A đến B là: $t=\dfrac{AB}{v_{CB}}=\dfrac{18}{21,6}=\dfrac{5}{6}\ h$.
\item Khi đi ngược dòng thì: $v_{CB}=v_{CN}-v_{NB}=16,2-5,4=10,8\ km/h$.
\item Thời gian canô đi từ B về A là: $t=\dfrac{AB}{v_{CB}}=\dfrac{18}{10,8}=\dfrac{5}{3}\ h$.
\item Thời gian canô đi từ A đến B rồi trở lại A là: $\dfrac{5}{6}+\dfrac{5}{3}=2,5\ h=2h30$ phút.
\end{itemize}	
}
\end{ex}

\begin{ex}%book %[1P1K9-1][Loai1] 
Hai bến tàu $M$ và $N$ cách nhau $60\ km$. Một tàu thủy đi xuôi dòng từ $M$ về $N$. Tàu thủy nghỉ lại ở $N$ trong $1\ h$ để bốc xếp hàng hóa rồi đi ngược dòng từ $N$ về lại $M$. Tổng cộng thời gian hết $10\ h$. Biết nước sông chảy với vận tốc $5\  km/h$. Vận tốc tàu thủy đi trong nước lặng là
\choice
{$20\ km/h$}
{$12\ km/h$}
{\True $15\ km/h$}
{$18,3\ km/h$}
\haidong{13}
\loigiai{	      
\begin{itemize}
\item Gọi v là tốc độ của tàu thủy trong nước lặng. 
\item Khi đi xuôi dòng, vận tốc của tàu thủy so với bờ là $v + 5$   km/h. 
\item Khi đi ngược dòng, vận tốc của tàu thủy so với bờ là $v - 5$   km/h.
\item Ta có phương trình $10 - 1 = 9 = \dfrac{60}{v + 5} + \dfrac{60}{v - 5}$ $\Rightarrow$ $3v^2 - 40v - 75 = 0$ $\Rightarrow v = 15$ km/h.
\end{itemize}	}
\end{ex}



%\loai{Tính thời gian vật chuyển động với vận tốc kéo theo}
\begin{ex}%book %[1P1G9-1][Loai1]
Một chiếc phà chạy xuôi dòng từ $A$ đến $B$ mất $3\ h$, khi chuyển động chạy ngược dòng về từ $B$ đến $A$ thì mất $6\ h$. Nếu phà tắt máy trôi theo dòng từ $A$ đến $B$ thì mất bao lâu?
\choice
{$6\ h$}
{\True $12\ h$}
{$7\ h$}
{$15\ h$}	
\haidong{13}
\loigiai{
\begin{itemize}
\item Ta có khi phà chạy xuôi dòng thì: $v_{PB}=v_{PN}+v_{NB} \Rightarrow AB=\left(v_{PN}+v_{NB}\right)t=\left(v_{PN}+v_{NB}\right).3$.
\item Khi phà chạy ngược dòng thì: $v_{PB}=v_{PN}-v_{NB} \Rightarrow AB=\left(v_{PN}-v_{NB}\right)t=\left(v_{PN}-v_{NB}\right).6 \Rightarrow \left(v_{PN}+v_{NB}\right).3=\left(v_{PN}-v_{NB}\right).6 \Leftrightarrow v_{PN}=3v_{NB}$.
\item Khi phà tắt máy thì vận tốc của phà so với bờ sông bằng vận tốc của nước so với bờ $\Rightarrow$ Khi tắt máy thì thời gian đi từ A - B là: $t=\dfrac{AB}{v_{NB}}=\dfrac{12v_{NB}}{v_{NB}}=12\ h$. 
\end{itemize}	
}
\end{ex}












%mac
%\loai{Tính vận tốc tương đối}

%\loai{Tính vận tốc tuyệt đối}
\begin{ex}%book %[1P1-9.2K][Loai1] %BAI 7+8 - HOC MAI
\immini[thm]{Một ca nô xuất phát từ điểm $A$ bên này sông, sang điểm $B$ bên kia sông theo phương vuông góc với bờ sông (hình vẽ). Vì nước chảy đều với vận tốc $3\ m/s$ nên ca nô đến bên kia sông tại điểm $C$ với vận tốc $5\ m/s$. Ca nô có vận tốc tuyệt đối là
\choice
{$2\ m/s$}
{$3\ m/s$}
{\True $4\ m/s$}
{$5\ m/s$}}{\begin{tikzpicture}[draw=Mapcolor, very thick,scale=1,thick,>=stealth']
\begin{scope}
\clip (2,0) -- (2,3) -- (8,3)  -- (8,0) -- cycle;
\foreach \x in {-2.9,-2.7,...,4}
{\draw[thin, dash pattern = on 4pt off 4pt] ([xshift = 4pt] + 2cm, \x cm) -- ++([xshift = -4pt]6.0cm, 0);}
\foreach \x in {-3.0,-2.8,...,4}
{\draw[thin, dash pattern = on 4pt off 4pt] ([xshift = 4pt] + 1.9cm, \x cm) -- ++([xshift = -4pt]6.0cm, 0);}
\draw (2,0) -- ++(0:6)(2,3) -- ++(0:6);  
\draw[dashed](4,0) -- ++(90:3)(4,0)--(7,3); 
\end{scope}
\node[rotate=90] at (4,0.3){\includegraphics[width=1cm]{images/boat}};
\node[below=3pt] at (4,0){$A$};
\node[above] at (4,3){$B$};
\node[above] at (7,3){$C$};
\end{tikzpicture}}
\haidong{13}	
\loigiai{
\begin{itemize}
\item Gọi $v_{c-b}$ là vận tốc của ca nô so với bờ sông.
\item $v_{c-n}$ là vận tốc của ca nô so với dòng nước: $v_{c-n} = 5\ m/s$
\item $v_{n-b}$ là vận tốc của nước so với bờ sông: $v_{n-b} = 3$ m/s
\item 	Ta có: $v^2_{c-n} = v^2_{c-b} + v^2_{n-b} \Rightarrow v_{c-b} = 4\ m/s$.
\end{itemize}	}
\end{ex}
\begin{ex}%Tailieuchuan
Một chiếc thuyền chạy xuôi dòng mất $2$ giờ, khi chạy về ngược dòng với cùng quãng đường thì mất $4$ giờ. Nếu thuyền tắt máy trôi theo dòng mất thời gian là
\choice
{$6$ giờ}
{$7$ giờ}
{\True $8$ giờ}
{$9$ giờ}
\loigiai{}
\end{ex}
\begin{ex}%Tailieuchuan
\immini[thm]{Một chiếc thuyền chuyển động từ điểm $A$ của bờ bên này đến điểm $B$ của bờ bên kia của con sông, do nước chảy siết thuyền không đến được bờ $B$ mà đến điểm $C$ cách bờ $B$ $180\ m$, biết sông rộng $240\ m$ và thời gian qua sông là $1$ phút. Vận tốc của thuyền so với dòng nước là
\choice
{$3 \ m/s$}
{\True $4 \ m/s$}
{$5 \ m/s$}
{$6 \ m/s$}}{	\begin{tikzpicture}[scale=1,thick,>=stealth',draw=Mapcolor]
\path
(3.5,0)coordinate(A)
(3.5,3)coordinate(B)
(1,3)coordinate(C)
;

\draw[Mapcolor] (0,0) -- (5,0);
\draw[Mapcolor] (0,3) -- (5,3);
\draw[Mapcolor] (A)--(B) (A)--(C) (B)--(C);
\node at (2,3.5) {$180\ m$};
\node at (4.5,1.5) {$240\ m$};

\foreach \x/\g in {A/-90,B/90,C/90}
\draw[fill=white](\x) circle(0.05)+(\g:0.3)node{$\x$};
\end{tikzpicture}}
\loigiai{
}
\end{ex}
\begin{ex}%Tailieuchuan
Một ô tô chạy với vận tốc $50\ km/h$ trong trời mưa. Mưa rơi theo phương thẳng đứng, trên cửa kính của xe, các vệt mưa rơi làm với phương thẳng đứng một góc $60^\circ$. Vận tốc của giọt mưa đối với xe ô tô là
\choice
{\True $57,73 \ km/h$}
{$50,36 \ km/h$}
{$45,45 \ km/h$}
{$60,50 \ km/h$}
\loigiai{

}
\end{ex}
\begin{ex}%Tailieuchuan
Một chiếc xe đang chạy với vận tốc $18 \ km/h$ trong mưa, giả sử mưa rơi thẳng đứng và đều đối với mặt đất. Người ngồi trên xe thấy các giọt mưa tạo một góc $30$ độ so với phương thẳng đứng. Tốc độ rơi của hạt mưa đối với mặt đất là
\choice
{$5 \ m/s$}
{\True $5 \sqrt{3} \ m/s$}
{$5 \ m/s$}
{$18 \sqrt{3} \ m/s$}

\loigiai{

}
\end{ex}
\begin{ex}%Tailieuchuan
Một chiếc thuyền chuyển động ngược dòng với vận tốc $14 \ km/h$ so với mặt nước, nước chảy với tốc độ $9 \ km/h$ so với bờ. Độ lớn vận tốc của thuyền so với bờ là
\choice
{$25 \ km/h$}
{$3,5 \ km/h$}
{\True $5 \ km/h$}
{$6 \ km/h$}
\loigiai{

}
\end{ex}





%\loai{Bài toán tổng hợp}
\textit{Sử dụng thông tin sau cho Câu 17 và Câu 18:} Một ca nô muốn đi thẳng qua một con sông rộng $0,10\ km$. Động cơ của ca nô tạo cho nó vận tốc $5,0\ km/h$ trong nước sông không chảy. Tuy nhiên, có một dòng chảy mạnh đang di chuyển về phía hạ lưu với vận tốc $3,0\ km/h$.

\begin{ex}%book %[1P1K9-1][Loai1] 
Ca nô phải đi theo hướng nào để đến vị trí ở bờ bên kia đối diện với vị trí xuất phát?
\choice
{$53,13^\circ$ xuôi theo hướng thượng lưu so với bờ sông}
{$53,13^\circ$ xuôi theo hướng hạ lưu so với bờ sông}
{\True $36,5^\circ$ xuôi theo hướng thượng lưu so với bờ sông}
{$36,5^\circ$ xuôi theo hướng hạ lưu so với bờ sông}
\loigiai{}
\end{ex}
\begin{ex}%book %[1P1K9-1][Loai1] 
Chuyến đi sẽ mất thời gian là 
\choice
{\True $1,5$ phút}
{$0,025$ phút}
{$15$ phút}
{$2,5$ phút}
\loigiai{
\immini{\begin{itemize}
\item[a)] $ \sin \alpha=\dfrac{v_n}{v_t}=\dfrac{3}{5}
\Rightarrow \alpha=36,87^\circ$\\
$\Rightarrow 90^{\circ}-36,87^{\circ}=53,13^{\circ}$, theo hướng thượng lưu so với bờ sông.
\item[b)] Chuyến đi sẽ mất:\\
$t=\dfrac{S}{v}=\dfrac{0,10}{\sqrt{5^2-3^2}}=0,025$ giờ $=1,5$ phút
\end{itemize}}{ \begin{tikzpicture}[scale=1,thick,>=stealth']
\tkzDefPoints{0/0/o, -1.5/1.5/a, 0/1.5/b}
\tkzDrawPoints[size=3,fill=white](o)
\draw[->] (0,0) -- (1.5,0);
\draw[->] (0,0) -- (-1.5,1.5);
\draw[->] (0,0) -- (0,1.5);
\tkzLabelAngles[pos=.7](a,o,b){$\alpha$}
\tkzMarkAngles[size=0.4cm](b,o,a)
\draw[dashed,thin] (-1.5,1.5) -- (0,1.5) -- (1.5,0);

\node at (-1,0.5) {$\vec v_t$};
\node at (0.2,0.7) {$\vec v$};
\node at (0.8,-0.25) {$\vec v_n$};
\draw[->] (-1.5,-0.5) -- (2.5,-0.5);
\node at (0.2,-0.9) {Dòng chảy};
\end{tikzpicture}}
}
\end{ex}













\subsubsection*{PHẦN 2. Thí sinh trả lời từ câu 1 đến câu 4. Trong mỗi ý a), b), c), d) ở mỗi câu, thí sinh chọn đúng hoặc sai.}
\setcounter{ex}{0}
\begin{ex}%Tailieuchuan
Một ca nô chạy qua sông xuất phát từ $A$, mũi ca nô hướng tới điểm $B$ ở bờ bên kia $AB$ vuông góc với bờ sông. Nhưng do nước chảy nên khi đến bờ bên kia, ca nô lại ở cách $B$ đoạn $240\ m$. Thời gian qua sông là $120\ s$. Nếu người lái giữ cho mũi ca nô chếch $60^\circ$ so với bờ sông và mở máy chạy như trước thì ca nô chạy tới đúng điểm $B$.
\choiceTF[t]
{\True Tốc độ dòng nước là $2 \ m/s$} 
{Tốc độ của ca nô so với bờ là $4 \ m/s$} 
{\True Bề rộng của dòng sông là $276\ m$} 
{\True Thời gian ca nô qua sông lần sau là $240\ s$} 
\loigiai{ 
\begin{itemchoice}
\itemch
\itemch
\itemch
\itemch
\end{itemchoice}
}
\end{ex}
\begin{ex}%Tailieuchuan
Một chiếc thuyền đi từ bến $A$ đến bến $B$ cách nhau $6\ km$ rồi trở về. Biết rằng tốc độ của thuyền trong nước yên lặng là $5\ km/h$, tốc độ nước chảy là $1 \ km/h$.
\choiceTF[t]
{\True Khi thuyền đi xuôi dòng, tốc kế của thuyền chỉ $6\ km/h$} 
{\True Vận tốc của thuyền khi đi ngược dòng là $4\ km/h$} 
{Thời gian chuyển động của thuyền trong quá trình trên là $2$ giờ $50$ phút} 
{Một người thuỷ thủ chạy từ đầu tàu đến đuôi tàu với tốc độ không đổi. Khi thuyền đi xuôi dòng và ngược dòng, tốc độ của người này so với nước là như nhau} 
\loigiai{ 
\begin{itemchoice}
\itemch
\itemch
\itemch
\itemch
\end{itemchoice}
}
\end{ex}
\begin{ex}%Tailieuchuan
Một phi công muốn máy bay của mình bay về hướng Tây trong khi gió thổi về hướng Nam với tốc độ $50\ km/h$. Biết rằng khi không có gió, vận tốc của máy bay là $200 \ km/h$.
\choiceTF[t]
{Khi không có gió, máy bay mất ít hơn $1$ giờ để bay quãng đường $200\ km$} 
{\True Để máy bay bay về hướng Tây trong khi gió thổi về hướng Nam, phi công phải lái máy bay theo hướng Tây - Bắc} 
{\True Máy bay được điều khiển lệch khỏi hướng Tây một góc xấp xỉ $14,5^\circ$} 
{Độ lớn vận tốc của máy bay so với mặt đất là $190,65\ km/h$} 
\loigiai{ 
\begin{itemchoice}
\itemch
\itemch
\itemch
\itemch
\end{itemchoice}
}
\end{ex}
\begin{ex}
\immini[thm]{Một ca nô chạy ngang qua một dòng sông, xuất phát từ $A$, hướng mũi về $B$. Sau $100\ s$, ca nô cập bờ bên kia ở điểm $C$ cách $B$ $200\ m$. Nếu người lái hướng mũi ca nô theo hướng $AD$ và vẫn giữ tốc độ máy như cũ thì ca nô sẽ cập bờ bên kia tại đúng điểm $B$.}{\begin{tikzpicture}[scale=1,thick,>=stealth',draw=Mapcolor]
\path
(3.5,0)coordinate(A)
(3.5,3)coordinate(B)
(4.5,3)coordinate(C)
(2,3)coordinate(D)
;
\draw[Mapcolor] (0,0) -- (5,0);
\draw[Mapcolor] (0,3) -- (5,3);
\draw[Mapcolor] (A)--(B) (A)--(C) (B)--(C) (A)--(D);

\foreach \x/\g in {A/-90,B/90,C/90, D/90}
\draw[fill=white](\x) circle(0.05)+(\g:0.3)node{$\x$};
\draw (3.32,0.36) .. controls (3.1,0.4) and (2.99,0.19) .. (3.1,0);
\node at (2.7,0.5) {$60^\circ$};
\draw[->,Mapcolor] (0,2) -- (1,2);
\draw[->,Mapcolor] (4.5,1.5) -- (5,1.5);
\draw[->,Mapcolor] (1,1) -- (1.5,1);
\draw[->,Mapcolor] (4,0.5) -- (4.5,0.5);
\end{tikzpicture}}
\choiceTF[t]
{\True Ca nô tham gia đồng thời hai chuyển động: chuyển động so với nước và chuyển động do nước kéo đi}
{\True Vận tốc của dòng nước so với bờ sông là $2 \ m/s$}
{Vận tốc của ca nô so với dòng nước là $3 \ m/s$}
{\True Chiều rộng của sông là $400\ m$}
\loigiai{
\begin{itemchoice}
\itemch Ca nô tham gia đồng thời 2 chuyển động: Chuyển động so với nước và chuyển động do nước kéo đi.
\itemch Vận tốc của dòng nước so với bờ sông là $v_{23}=\dfrac{\boldsymbol{B C}}{\boldsymbol{t}}=\dfrac{\mathbf{2 0 0}}{\mathbf{1 0 0}}=\mathbf{2 \  m} /\ s$.
\itemch Gọi $\vec{v_{12}}$ là vận tốc của ca nô so với dòng nước.

$\vec{v_{23}}$ là vận tốc của dòng nước so với bờ sông.

$\vec{v_{13}}$ là vận tốc của ca nô so với bờ sông.

Từ hình vẽ ta có: $A D=2 \cdot D B$ nên $v_{12}=2 \cdot v_{23}=4 \ m/s$.
\itemch Chiều rộng của bờ sông là $A B=v_{12} \cdot t=4.100=400 \ m$.
\end{itemchoice}
}
\end{ex}


\subsubsection*{PHẦN 3. Thí sinh trả lời từ câu 1 đến câu 6.}
\setcounter{ex}{0}
\textit{Sử dụng thông tin sau cho Câu 1 và Câu 2:} Trong trận lũ lụt tại miền Trung vào tháng $10/2020$, dòng lũ có tốc độ lên đến khoảng $4\ m/s$. Bộ Quốc Phòng đã trang bị ca nô công suất lớn trong công tác cứu hộ. Trong một lần cứu hộ, đội cứu hộ đã sử dụng ca nô chạy với tốc độ $8\ m/s$ so với dòng nước để cứu những người gặp nạn đang mắc kẹt trên một mái nhà cách trạm cứu hộ khoảng $2\ km$.Biết đội cứu hộ phải đi xuôi dòng lũ.
\begin{ex}%book
Sau bao nhiêu giây đội cứu hộ đến được chỗ người bị nạn (làm tròn kết quả đến chữ số hàng đơn vị)?
\shortans[oly]{167} 
\loigiai{
\begin{itemize}
\item $v_{13}$ là vận tốc tuyệt đối của ca nô so với bờ
\item $v_{12}$ là vận tốc tương đối của ca nô so với dòng nước
\item $v_{23}$ là vận tốc kéo theo của dòng nước so với bờ
\item Công thức cộng vận tốc: $\vec v_{13}=\vec v_{12}+\vec v_{23}$
\item[a)] Đội cứu hộ đi xuôi dòng nên $v_{13}=v_{12}+v_{23}=8+4=12\ m/s$.
\item Thời gian đội cứu hộ đến được chỗ bị nạn: $t=\dfrac{d}{v_{13}}=\dfrac{2000}{12}=\dfrac{500}{3}\ s=$ 2 phút 46 giây
\end{itemize}
}
\end{ex}
\begin{ex}%book
Sau khi cứu người, đội cứu hộ phải mất bao nhiêu giây để quay lại trạm ban đầu?
\shortans[oly]{500}
\loigiai{
\begin{itemize}
\item $v_{13}$ là vận tốc tuyệt đối của ca nô so với bờ
\item $v_{12}$ là vận tốc tương đối của ca nô so với dòng nước
\item $v_{23}$ là vận tốc kéo theo của dòng nước so với bờ
\item Công thức cộng vận tốc: $\vec v_{13}=\vec v_{12}+\vec v_{23}$
\item[a)] Đội cứu hộ đi xuôi dòng nên $v_{13}=v_{12}+v_{23}=8+4=12\ m/s$.
\item Thời gian đội cứu hộ đến được chỗ bị nạn: $t=\dfrac{d}{v_{13}}=\dfrac{2000}{12}=\dfrac{500}{3}\ s=$ 2 phút 46 giây
\item[b)] Khi quay trở lại vị trí ban đầu thì đội cứu hộ đi ngược dòng nước nên: $v'_{13}=v_{12}-v_{23}=8-4=4\ m/s$.
\item Thời gian để đội cứu hộ trở lại vị trí ban đầu: $t=\dfrac{d}{v'_{13}}=\dfrac{2000}{4}=500\ s=$8 phút 20 giây
\end{itemize}
}
\haidong{15}
\end{ex}


\textit{Sử dụng thông tin sau cho Câu 3, Câu 4 và Câu 5:} 	Một người đang ở phía Tây của một cái hồ và muốn bơi ngang qua để đến vị trí ở phía Đông, đối diện với vị trí xuất phát của mình. Người này có thể bơi với vận tốc $1,9\ m/s$ khi nước hồ lặng. Biết rằng lá cây trôi trên mặt nước hồ được $4,2\ m$ về hướng nam trong $5\ s$.
\begin{ex}%book %[1P1K9-1][Loai1]
Người này sẽ phải bơi theo hướng từ Đông sang Bắc bao nhiêu độ để đến vị trí đối diện trực tiếp với vị trí của anh ta (làm tròn kết quả đến chữ số hàng phần mười)? 
\shortans[oly]{23,5}
\loigiai{\begin{itemize}
\item[a)] Vận tốc dòng nước: $v_n=\dfrac{42}{3,5}=0,84\ m/s$.\\
$\tan\alpha=\dfrac{v_n}{v_{ng}}=\dfrac{0,84}{1,9}\Rightarrow \alpha=24,9^\circ$.\\
$v=\sqrt{v_n^2+v_{ng}^2}=2,1\ m/s$\\
Vậy để người đó bơi sang đúng vị trí đối diện với vị trí xuất phát thì người đó phải bơi với vận tốc 2,1 m/s theo hướng Bắc hợp với hướng Đông với góc $23,5^\circ$.
\end{itemize}}
\end{ex}

\begin{ex}
Vận tốc tổng hợp của người đó có độ lớn là bao nhiêu $m/s$?
\shortans[oly]{1,9}
\loigiai{
Vận tốc tổng hợp: $v_{ng} = 1,9\ m/s$ theo hướng Đông.}
\end{ex}
\begin{ex}
Nếu hồ rộng $4,8\ km$ thì người đó phải bơi bao nhiêu phút (làm tròn kết quả đến chữ số hàng đơn vị)?
\shortans[oly]{42}
\loigiai{Thời gian để bơi hết hồ: $t=\dfrac{d}{v_{ng}}=\dfrac{4800}{1,9}=2526\ s=42$ phút.}
\end{ex}
\begin{ex}%Tailieuchuan
Một xe đạp chuyển động thẳng đều với tốc độ lúc không có gió là $15\ km/h$. Người này đi từ $A$ tới $B$ xuôi gió và đi từ $B$ trở lại $A$ ngược gió. Biết vận tốc gió là $1\ km/h$ và $A$ cách $B$ $28\ km$. Tổng thời gian đi và về mất bao nhiêu giờ?
\shortans[oly]{3,75}
\loigiai{
}
\end{ex}







